\section{Information: Selection, Signaling \& Agency}

\subsection{Adverse Selection (Akerlof)}
Quality $\theta \sim F$ known to seller only; buyers pay expected quality. Price $P$ attracts only sellers with $\theta \leq P$ $\implies$ average quality traded is $\E[\theta \mid \theta \leq P] < P$ for valuable cars $\implies$ demand shifts down $\implies$ \textbf{unraveling}: possibly no trade at all, even when every trade has positive surplus.

\textbf{Key condition:} selection is on something \textit{priced}. If buyers value quality and sellers know it, the marginal seller drags down the average.

\textit{Application (insurance):} sickest buy the most coverage $\implies$ premiums reflect sick pool $\implies$ healthy drop out (death spiral). Mandates or group pooling (employer plans) fix the margin.

\textit{TFU trap:} adverse selection is about \textbf{hidden type} (pre-contract); moral hazard is \textbf{hidden action} (post-contract). Both can shrink the same market; identify which margin the question moves.

\subsection{Market Responses to Selection}
No regulator needed for partial fixes---the market prices information:
\begin{itemize}[nosep]
  \item \textbf{Warranties/guarantees}: cheap for high quality to offer (low expected repair cost) $\implies$ separates types
  \item \textbf{Reputation \& brand}: repeat play converts hidden type into observable history
  \item \textbf{Certification/inspection}: third parties sell the information directly (CARFAX, ratings)
  \item \textbf{Screening}: uninformed side offers a menu (deductibles); risk types self-select
\end{itemize}

\subsection{Signaling (Spence)}
Informed side burns resources to reveal type. Education costs $c/\theta$ per unit (cheaper for high ability $\theta$); employers pay expected productivity.

\textbf{Single crossing:} signal marginal cost decreasing in type $\implies$ separating equilibrium exists where $s^*(\theta)$ increasing and wages $= \theta$.

\textbf{Separation condition:} low type prefers not to mimic: $w_H - c\,s_H/\theta_L \leq w_L$.

\textbf{Welfare:} signaling is \textit{privately} rational but can be \textit{socially} wasteful---output unchanged, resources spent separating. Contrast with human capital view where education raises productivity; empirically both operate.

\textit{TFU trap:} a subsidy to education can \textit{reduce} welfare in a pure signaling model (cheaper signal $\implies$ more of it needed to separate).

\subsection{Principal--Agent / Moral Hazard}
Principal offers wage contract $w(\cdot)$; agent chooses hidden effort $e$ at cost $c(e)$; output signal $y$ depends stochastically on $e$.

\textbf{First best} (observable $e$): full insurance for risk-averse agent, wage flat, effort dictated.

\textbf{Second best:} incentive compatibility (IC) forces wage to vary with $y$ $\implies$ risk-averse agent bears risk $\implies$ \textbf{risk--incentives trade-off}. Power of incentives increases with: agent risk tolerance, informativeness of $y$ about $e$, effort responsiveness.

\textbf{Informativeness principle:} pay conditions on any signal informative about effort (relative performance evaluation filters common shocks).

\textbf{Monitoring vs.\ punishment:} with costly monitoring (probability $p$ of detection), large penalties economize on enforcement: expected penalty $pf$ deters, so raise $f$, lower $p$ (mirrors Becker crime logic). \textit{Core 1996 (Exxon):} punitive damages $10\times$ actual harm are efficient if detection probability is $1/10$.

\textit{Applications:} sharecropping (output share balances risk vs.\ incentives); efficiency wages (pay above market so job loss is a real penalty); deductibles in insurance.
